\documentclass[8pt]{extarticle}
\usepackage[margin=0.45in]{geometry}
\usepackage{graphicx,booktabs,amsmath,microtype,adjustbox,multicol}
\usepackage[table]{xcolor}
\usepackage{caption}
\captionsetup{font=footnotesize,labelfont=bf}
\usepackage{titlesec}
\titlespacing{\section}{0pt}{4pt}{2pt}
\titleformat{\section}{\bfseries}{\thesection.}{0.5em}{}
\pagestyle{empty}
\setlength{\parskip}{2pt}\setlength{\parindent}{0pt}
\graphicspath{{../figures/}}

\begin{document}

\begin{center}
{\Large\bfseries Accuracy is not mechanism: a causal roll-audit of PPG blood-pressure models
and their out-of-distribution failure}\\[2pt]
{\footnotesize Mechanistic-faithfulness audit across VitalDB, MIMIC-BP, and four external PPG datasets}
\end{center}

\vspace{-4pt}
\noindent\rule{\textwidth}{0.4pt}

\textbf{Abstract.} Deep photoplethysmography (PPG) blood-pressure (BP) models are compared on
in-distribution error, but error alone cannot say \emph{whether a model uses the physiology it
claims}. We train five architectures on VitalDB and audit them with (i) a causal \emph{roll
test} that shifts the PPG channel to lengthen pulse arrival time (PTT) and measures the sign of
the BP response, and (ii) a 31-cue linear-probe battery read across layers. All five models
reach near-identical in-distribution DBP error ($\sim$8~mmHg) yet differ up to $8\times$ in how
strongly and correctly they use PTT, and this mechanism gap---not the in-distribution
error---predicts a large out-of-distribution (OOD) collapse ($+2$ to $+13$~mmHg) across five
external datasets. Distilling the single cue that survives the audit into a gradient-boosted
tree matches deep in-distribution accuracy and generalizes better; adding age and sex lowers
error by $0.88$~mmHg on a clinical cohort.

\vspace{1pt}
\begin{minipage}[t]{0.53\textwidth}
\vspace{-2pt}
\includegraphics[width=\linewidth]{fig0_design.png}\\[-1pt]
{\scriptsize\textbf{Figure~1.} (a)~Design. (b)~Example ECG/PPG waveforms: the morphology shift the
model faces. (c)~SBP/DBP distribution shift. (d)~The causal roll-audit.}
\end{minipage}\hfill
\begin{minipage}[t]{0.44\textwidth}
\vspace{-2pt}
\captionof{table}{Datasets and measured distribution shift.}
\adjustbox{max width=\linewidth}{\scriptsize\begin{tabular}{lrrccccc}
\toprule
Dataset & Segments & Subjects & SBP / DBP (mmHg) & Age & Channels & Demographics & KS \\
\midrule
VitalDB-Vital (train) & 465,480 & 1,293 & 115$\pm$19 / 63$\pm$12 & 59 & ECG+PPG & age, sex, BMI & -- \\
VitalDB-Vital (CalFree test) & 57,600 & 144 & 115$\pm$19 / 63$\pm$12 & 58 & ECG+PPG & age, sex, BMI & -- \\
MIMIC-BP & 137,160 & 1,524 & 113$\pm$18 / 58$\pm$12 & -- & ECG+PPG & -- & 0.45 \\
BCG & 3,063 & 40 & 121$\pm$15 / 67$\pm$9 & 34 & PPG & age, sex & 0.54 \\
Sensors & 11,102 & 1,195 & 134$\pm$22 / 65$\pm$11 & 57 & PPG & age, sex & 0.44 \\
UCI2 & 410,596 & 10,793 & 132$\pm$22 / 67$\pm$11 & -- & PPG & -- & 0.49 \\
PPG-BP & 619 & 218 & 128$\pm$20 / 72$\pm$11 & 57 & PPG & age, sex, BMI & 0.52 \\
\bottomrule
\end{tabular}}
\vspace{3pt}
\captionof{table}{ID vs.\ OOD DBP error and audited mechanism.}
\adjustbox{max width=\linewidth}{\scriptsize\begin{tabular}{lrrrrrrc}
\toprule
& \multicolumn{6}{c}{DBP MAE (mmHg)} & Mechanism \\
\cmidrule(lr){2-7}
Model & ID & MIMIC & BCG & Sensors & UCI2 & PPG-BP & PAT slope / \%correct \\
\midrule
lenet1d & 8.6 & 13.1 & 11.0 & 9.5 & 9.7 & 11.1 & -15 / 68\% \\
inception1d & 8.0 & 13.5 & 15.2 & 13.2 & 13.4 & 17.0 & -25 / 68\% \\
xresnet1d50 & 8.4 & 15.1 & 18.0 & 15.1 & 14.5 & 20.5 & -40 / 77\% \\
xresnet1d101 & 8.2 & 17.1 & 18.1 & 15.9 & 15.7 & 21.4 & -33 / 77\% \\
transformer & 8.2 & 10.7 & 13.0 & 11.9 & 11.9 & 16.2 & -5 / 54\% \\
\bottomrule
\end{tabular}}
\vspace{3pt}
\captionof{table}{Interpretable model from audit-passing features.}
\adjustbox{max width=\linewidth}{\scriptsize\begin{tabular}{lr}
\toprule
Quantity & Value \\
\midrule
Faithful model (feature source) & xresnet1d50 \\
Audit-passing features & pat \\
ID DBP MAE (mmHg) & 8.89 \\
OOD DBP MAE -- ppgbp & 11.18 \\
OOD DBP MAE -- sensors & 8.05 \\
OOD DBP MAE -- bcg & 8.07 \\
OOD DBP MAE -- uci2 & 8.63 \\
OOD DBP MAE -- mimic-bp & 9.87 \\
\midrule
Age/sex ablation set & ppgbp \\
DBP MAE with age+sex & 8.84 \\
DBP MAE without & 9.72 \\
Improvement from demographics & +0.88 mmHg \\
SHAP mean$|\cdot|$: age / sex & 1.97 / 0.58 \\
\bottomrule
\end{tabular}}
\end{minipage}

\vspace{4pt}
\begin{multicols}{2}
\section{Setup}
Train: VitalDB, subject-disjoint splits. OOD: MIMIC-BP (ECG+PPG, 1{,}524 patients---the
roll-audit runs here) plus BCG, Sensors, UCI2, PPG-BP (PPG-only). Distribution shift is measured
per cue (KS vs.\ VitalDB, Table~1), so ``OOD'' is quantified rather than assumed. External
accuracy is reported with subject-level bootstrap CIs.

\section{Key findings}
\textbf{(1) Same accuracy, different fate.} In-distribution DBP error is flat across models
($8.0$--$8.6$~mmHg); OOD error fans out to $21$~mmHg. \emph{Larger models generalize worse}:
XResNet101 is best in-distribution but worst on every OOD set; the smallest model (LeNet)
transfers best.\\
\textbf{(2) Mechanism predicts the gap.} The causal PAT slope (Table~2) ranges from $-40$
(strongly, correctly physiological) to $-5$~mmHg/s (barely used) at identical in-distribution
error, and tracks OOD fragility.\\
\textbf{(3) Decodable $\neq$ used.} Many cues are linearly decodable (e.g.\ pulse width
$R^2\!=\!0.70$) yet fail the causal sign test---the faithfulness signature. For the strongest
model only \emph{PTT} passes both.\\
\textbf{(4) Faithful features transfer.} A tree built from the audit-passing cue matches deep
in-distribution accuracy and beats every deep model on MIMIC-BP; age/sex cut error by
$0.88$~mmHg on PPG-BP (SHAP: age $1.97$, sex $0.58$).
\end{multicols}

\vspace{2pt}
\noindent\rule{\textwidth}{0.4pt}
{\footnotesize\textbf{Takeaway.} An in-distribution leaderboard hides which models are right for
the right reason. A causal input-space audit exposes it, predicts OOD failure, and points to the
features an interpretable, better-generalizing model should use.}

\end{document}
